\section{Method}
\label{sec:method}

Fig.~\ref{fig:pipeline} shows the pipeline. Each passage is indexed twice, in
BM25 and in a dense vector index. A query goes to both retrievers, their top-100
lists are merged by weighted Reciprocal Rank Fusion (RRF) with a weight that is
either constant or, in our method, predicted per query, and a cross-encoder
reorders the head of the fused list before the top five passages reach the
generator.

\begin{figure*}[t]
\centering
\begin{tikzpicture}[
  font=\footnotesize,
  box/.style={draw, rounded corners=2pt, align=center, minimum height=9mm, inner sep=3pt, fill=gray!8},
  ours/.style={box, fill=blue!10, draw=blue!60!black, thick},
  store/.style={box, fill=white, dashed},
  arr/.style={-{Stealth[length=2mm]}, semithick}]
  % offline row
  \node[box]   (docs)  at (-3.1,2.0) {Documents};
  \node[box]   (chunk) at (-0.2,2.0) {Structural chunking\\(300 words, 50 overlap)};
  \node[store] (idx)   at (3.4,2.0)  {BM25 index (bm25s)\\HNSW index (FAISS)};
  \draw[arr] (docs) -- (chunk);
  \draw[arr] (chunk) -- (idx);
  % query row
  \node[box]  (q)    at (0,0)    {Query $q$\\(optional LLM expansion)};
  \node[box]  (lists) at (3.4,0) {Sparse top-100\\Dense top-100};
  \coordinate (join) at (5.2,0);
  \node[ours] (qaf)  at (6.4,-1.5) {QAF: 12 features\\ridge $\rightarrow \alpha(q)$};
  \node[box]  (rrf)  at (7.6,0)  {Weighted RRF\\top-100};
  \node[box]  (rr)   at (10.2,0) {Cross-encoder\\rerank top-50};
  \node[box]  (gen)  at (13.0,0) {Qwen2.5-7B\\top-5, cited answer};
  \draw[arr] (q) -- (lists);
  \draw[arr, dotted] (idx) -- (lists);
  \draw[semithick] (lists.east) -- (join);
  \draw[arr] (join) -- (rrf.west);
  \draw[arr] (join) |- (qaf.west);
  \draw[arr] (qaf.east) -| node[pos=0.75, right]{$\alpha$} (rrf.south);
  \draw[arr] (rrf) -- (rr);
  \draw[arr] (rr) -- (gen);
\end{tikzpicture}
\caption{Pipeline. Top: offline indexing. Bottom: query time. The shaded block
is Query-Adaptive Fusion (QAF), which sets the dense weight $\alpha$ of the
fusion step separately for every query.}
\label{fig:pipeline}
\end{figure*}

\subsection{Chunking}
\label{sec:chunking}
BEIR and HotpotQA passages are already short and are indexed as they are.
Chunking matters for the RFC corpus, where one document can run past 100
pages. Our chunker splits the cleaned text into blocks at blank lines. A
one-line block is a heading if it is a numbered title of at most 12 words that
does not end in a period, colon or comma (``4.2.1 Frame Header''), or a short
line of at most 8 words that starts with a capital and contains only letters
and spaces. A heading closes the current chunk and becomes the first line of
the next, so section names are indexed with their bodies. Paragraphs are added
until the chunk would pass 300 words (hard ceiling 380); longer paragraphs are
split at sentences, then into word windows. Consecutive chunks of the same
section share 50 words of overlap. Each chunk is a verbatim slice of its source
and stores its character offsets, which we use for relevance labels. An
optional break on low embedding similarity between paragraphs was off in every
reported run, so boundaries follow document structure only.

\subsection{Sparse retrieval}
BM25 scores the title and text of each passage with the Lucene variant
implemented in \texttt{bm25s}~\cite{lu2024bm25s}:
\begin{equation}
  s_{\mathrm{BM25}}(q,d)=\sum_{t\in q}\mathrm{idf}(t)\,
  \frac{f(t,d)}{f(t,d)+k_1\!\left(1-b+b\,\frac{|d|}{\mathrm{avgdl}}\right)},
  \label{eq:bm25}
\end{equation}
where $f(t,d)$ is the frequency of $t$ in $d$ and
$\mathrm{idf}(t)=\ln\bigl(1+\frac{N-n_t+0.5}{n_t+0.5}\bigr)$ for $N$ passages,
$n_t$ of which contain $t$~\cite{robertson2009bm25}; $k_1=1.2$, $b=0.75$.

The tokenizer is where we departed most from defaults. The default
\texttt{bm25s} tokenizer keeps only runs of two or more word characters, so
\texttt{HTTP/2} becomes \texttt{http} and the version number, often exactly
what an RFC question asks about, is lost. Ours lowercases the text and
first matches compounds, alphanumeric runs joined by \texttt{.}, \texttt{\_},
\texttt{/}, \texttt{:} or \texttt{-}, and emits each compound both whole and as
its parts (\texttt{http/2} gives \texttt{http/2}, \texttt{http}, \texttt{2}).
Purely alphabetic tokens go through an English stopword list and the Snowball
stemmer; tokens with a digit or separator are never stemmed or dropped. The
same function serves documents, queries and the IDF look-ups used by QAF.

\subsection{Dense retrieval}
We compare two bi-encoders: BGE-small-en-v1.5~\cite{xiao2024cpack,baai2023bgesmall}
(384 dimensions) and EmbeddingGemma~2~\cite{vera2026embeddinggemma2}, the
successor of EmbeddingGemma~\cite{vera2025embeddinggemma}, whose 768-dimensional
vectors are trained with a Matryoshka objective~\cite{kusupati2022matryoshka}
so that prefixes of 512, 256 or 128 dimensions remain usable. We load only the
text tower of EmbeddingGemma~2, in 32-bit floats. Each model gets its
documented prompts: BGE prefixes queries with ``Represent this sentence for
searching relevant passages:'' and encodes passages as plain text, while
EmbeddingGemma~2 uses ``task: search result \textbar{} query:'' for queries
and ``title: \{title or none\} \textbar{} text: \{passage\}'' for passages.
Vectors are L2-normalised and stored in a FAISS~\cite{johnson2021faiss} HNSW
graph~\cite{malkov2020hnsw} ($M=32$, $\mathit{efConstruction}=200$,
$\mathit{efSearch}=128$) searched by inner product.

\subsection{Weighted reciprocal rank fusion}
BM25 scores and cosine similarities are on different scales, so we fuse
ranks~\cite{cormack2009rrf}. For a dense weight $\alpha\in[0,1]$,
\begin{equation}
  s_{\alpha}(d)=\frac{\alpha}{k+r_{\mathrm{dense}}(d)}
               +\frac{1-\alpha}{k+r_{\mathrm{sparse}}(d)},\qquad k=60,
  \label{eq:rrf}
\end{equation}
where $r(\cdot)$ is the 1-based rank in each top-100 list and a passage absent
from a list gets nothing from it. The fused list keeps 100 passages, with ties
broken by passage id. At $\alpha=0.5$ this is plain RRF up to a constant
factor; $\alpha=1$ and $\alpha=0$ give back the dense and sparse rankings.

\subsection{Query-Adaptive Fusion (QAF)}
\label{sec:qaf}
One $\alpha$ for all queries is a compromise. ``Which opcode marks a WebSocket
close frame?'' needs the sparse side; a loosely worded finance question needs
the dense side. QAF predicts $\alpha$ per query from 12 features
(Table~\ref{tab:features}) that reuse work the pipeline has already done.
Fig.~\ref{fig:qaf} shows how the features are grouped, how $\alpha$ is
produced, and how the model is trained.

\begin{figure}[t]
\centering
\begin{tikzpicture}[
  font=\footnotesize,
  box/.style={draw, rounded corners=2pt, align=center, inner sep=2.5pt, fill=gray!8},
  feat/.style={box, fill=white, text width=1.62cm, minimum height=1.55cm},
  train/.style={box, fill=white, text width=3.7cm},
  arr/.style={-{Stealth[length=1.6mm]}, semithick}]
  % inputs
  \node[box] (q)  at (0.95,0) {Query $q$};
  \node[box] (sp) at (3.05,0) {BM25 top-100};
  \node[box] (de) at (7.25,0) {Dense top-100};
  % feature groups from qaf.FEATURE_NAMES (agreement drawn between the two lists)
  \node[feat] (lex) at (0.95,-1.55) {\textbf{Lexical (7)}\\id, num, caps ratios; length; mean/max IDF; OOV};
  \node[feat] (sc)  at (3.05,-1.55) {\textbf{Sparse conf. (2)}\\top score, margin};
  \node[feat] (dc)  at (7.25,-1.55) {\textbf{Dense conf. (2)}\\top score, margin};
  \node[feat] (ag)  at (5.15,-1.55) {\textbf{Agreement (1)}\\overlap@10 of the two lists};
  \draw[arr] (q) -- (lex);
  \draw[arr] (sp) -- (sc);
  \draw[arr] (sp) -- (ag.north);
  \draw[arr] (de) -- (dc);
  \draw[arr] (de) -- (ag.north);
  % model
  \node[box, text width=3.7cm] (ridge) at (2.3,-3.2) {StandardScaler + Ridge\\(regularisation 1.0)};
  \node[box, text width=2.3cm] (clip)  at (6.6,-3.2) {$\alpha(q)$, clipped\\to $[0,1]$};
  % feature bus: all 12 features go into the scaler + ridge
  \foreach \n in {lex,sc,ag,dc} \draw[semithick] (\n.south) -- (\n.south |- 0,-2.55);
  \draw[semithick] (lex.south |- 0,-2.55) -- (dc.south |- 0,-2.55);
  \draw[arr] (ridge.north |- 0,-2.55) -- (ridge.north);
  \draw[arr] (ridge) -- (clip);
  \node[box, text width=2.3cm] (rrf) at (6.6,-4.35) {Weighted RRF\\(Eq.~\eqref{eq:rrf})};
  \node[box, text width=2.3cm] (rr)  at (6.6,-5.4)  {Cross-encoder\\rerank top-50};
  \draw[arr] (clip) -- (rrf);
  \draw[arr] (rrf) -- (rr);
  % training (offline)
  \node[train] (t3) at (2.3,-4.6) {Plateau-centre target $\alpha^{o}(q)$\\(0.5 if no signal)};
  \node[train] (t2) at (2.3,-5.6) {Sweep $\alpha\in\{0,0.1,\dots,1\}$:\\nDCG@10 of fused list};
  \node[train] (t1) at (2.3,-6.55) {SciFact-dev + FiQA-dev queries};
  \draw[arr] (t1) -- (t2);
  \draw[arr] (t2) -- (t3);
  \draw[arr] (t3) -- node[right, pos=0.4]{fit Ridge} (ridge);
  \node[draw, dashed, rounded corners=3pt, inner sep=3pt, fit=(t1)(t2)(t3)] (tr) {};
  \node[anchor=north west, inner sep=1pt] at (tr.south west) {\scriptsize training (offline)};
\end{tikzpicture}
\caption{QAF at query time (top and right) and its offline training (dashed
box). The four feature groups together hold the 12 features of
Table~\ref{tab:features}.}
\label{fig:qaf}
\end{figure}


\begin{table}[t]
\centering
\caption{QAF query features. The query is split into alphanumeric tokens
(compounds kept whole); ``ratio'' means the share of those tokens.}
\label{tab:features}
\footnotesize
\begin{tabular}{@{}>{\raggedright\arraybackslash}p{0.30\columnwidth}>{\raggedright\arraybackslash}p{0.64\columnwidth}@{}}
\toprule
Feature & Definition \\
\midrule
\texttt{id\_ratio} & identifier-like tokens: letters mixed with digits, an internal \texttt{\_-./}, or all capitals \\
\texttt{num\_ratio} & tokens made only of digits \\
\texttt{caps\_ratio} & all-capital tokens of length $\ge 2$ (acronyms) \\
\texttt{len\_tokens} & number of tokens \\
\texttt{mean\_idf}, \texttt{max\_idf} & BM25 IDF of the query terms in the target corpus \\
\texttt{oov\_ratio} & tokens absent from the dense encoder's tokenizer vocabulary (whole-word entries) \\
\texttt{sparse\_top}, \texttt{sparse\_margin} & top BM25 score and its gap to the second \\
\texttt{dense\_top}, \texttt{dense\_margin} & top cosine score and its gap to the second \\
\texttt{overlap\_at10} & share of the dense top-10 also in the sparse top-10 \\
\bottomrule
\end{tabular}
\end{table}

\textit{Training target.} For a labelled training query we sweep $\alpha$ over
$\{0,0.1,\dots,1\}$ and compute nDCG@10 of the fused list at each value. Since
Eq.~\eqref{eq:rrf} uses ranks only, nDCG is piecewise constant in $\alpha$ and
the optimum is usually a run of equally good grid points. Taking the first or
last maximiser would push targets toward the ends of the range for no real
reason, so the oracle $\alpha^{o}(q)$ is the centre of the longest contiguous
run of grid values that reach the best nDCG@10; between equally long runs we
keep the one nearer $0.5$. A query with nDCG@10 of zero at every $\alpha$
carries no signal and gets $0.5$.

\textit{Model and cost.} Standardised features feed a ridge regression
(regularisation $1.0$) whose output, clipped to $[0,1]$, is used as $\alpha$.
We fit one model per encoder on the pooled SciFact and FiQA dev queries; the
RFC and HotpotQA collections are never seen in training. We chose a linear
model so that its coefficients can be read and checked against intuition.
Every feature comes from the query, the two candidate lists, the BM25
document frequencies or, for \texttt{oov\_ratio}, the encoder's tokenizer
vocabulary, so QAF adds feature extraction and a 12-term dot product, with no
model call.

\subsection{Reranking, post-processing and generation}
The first 50 passages of the fused list are rescored by the BGE-reranker-base
cross-encoder, which reads query and passage (title and text, up to 512 tokens)
jointly as in monoBERT~\cite{nogueira2019passage}; the top five passages of the
new order form the context. The interactive pipeline can also drop near
duplicates (Jaccard similarity of word 5-gram sets $\ge 0.9$) and apply a
metadata filter before reranking; both are disabled in the reported
experiments. Optional query expansion, used only in the ablation, asks the
local LLM for at most 20 keywords or synonyms and appends them to the query.

Answers come from Qwen2.5-7B-Instruct~\cite{qwen2024qwen25} served by Ollama
at temperature 0 with a fixed seed. The prompt numbers the passages
\texttt{[1]}--\texttt{[5]} and asks for a short answer citing them in that
form, or the exact reply ``not in context'' when they do not support one.
Cited numbers are mapped back to passage ids, so we can check whether a cited
passage is relevant.
